\documentclass[twocolumn, a4paper, 10pt]{ctexart}

% ================= 包引入 =================
\usepackage{geometry}
\geometry{left=1.5cm, right=1.5cm, top=2.5cm, bottom=2.5cm} % 页边距
\usepackage{fancyhdr} % 页眉页脚
\usepackage{graphicx} % 图片
\usepackage{amsmath, amsfonts, amssymb} % 数学公式
\usepackage{booktabs} % 三线表
\usepackage{caption} % 图表标题
\usepackage{titlesec} % 标题格式
\usepackage{enumitem} % 列表格式
\usepackage{xcolor} % 颜色
\usepackage{float}
\usepackage{balance} % 最后一页双栏平衡
\usepackage{bm} % 粗体数学符号

% ================= 格式设置 =================
% 图表标题设置 (字体设为小五)
\captionsetup{font=small, labelsep=quad}
\renewcommand{\figurename}{图}
\renewcommand{\tablename}{表}

% 节标题设置
\titleformat{\section}{\zihao{4}\sffamily\bfseries}{\thesection}{1em}{}
\titleformat{\subsection}{\zihao{5}\sffamily\bfseries}{\thesubsection}{1em}{}
\titleformat{\subsubsection}{\zihao{-5}\sffamily\bfseries}{\thesubsubsection}{1em}{}

% 页眉页脚设置 (模仿第一篇论文)
\pagestyle{fancy}
\fancyhf{}
\fancyhead[L]{\small 第 8 卷第 2 期 \\ 2024 年 6 月}
\fancyhead[C]{\small \bfseries 物 联 网 学 报 \\ Chinese Journal on Internet of Things}
\fancyhead[R]{\small Vol.8 No.2 \\ June 2024}
\renewcommand{\headrulewidth}{0.5pt}

% 首行缩进
\setlength{\parindent}{2em}

% ================= 文档开始 =================
\begin{document}

% 使用 \twocolumn 的可选参数来实现单栏标题和摘要跨双栏
\twocolumn[
\begin{center}
    \vspace*{-1cm}
    % 中文标题 (根据英文标题为您补充)
    {\zihao{3} \sffamily \bfseries 基于热插拔混合专家的多视角结构化音乐生成控制机制\par}
    \vspace{4mm}
    
    % 中文作者
    {\zihao{4} 作者一$^1$，作者二$^2$，作者三$^1$\par}
    \vspace{2mm}
    
    % 中文机构
    {\zihao{-5} （1. XX大学 计算机学院，省份 城市 邮编；2. XX大学 人工智能研究院，省份 城市 邮编）\par}
    \vspace{6mm}
    
    % 中文摘要与关键词
    \parbox{0.9\textwidth}{
        \zihao{5}
        {\bfseries 摘 要：}{\textbf{(1) 问题：}}AI 音乐生成与视频生成之间存在显著的提示词工程代差——视频领域已通过标准化的镜头语言实现多维度精细控制，而音乐生成仍局限于粗粒度标签组合。{\textbf{(2) 方案：}}本文提出 Camera-Music ControlNet，将多视角叙事控制方法论系统性迁移至音乐生成领域。{\textbf{(3) 创新：}}核心创新包括建立九组跨模态结构性映射，定义覆盖六维的可控 Token 词汇表；构建基于显式 [STYLE] 路由的混合专家架构，通过 Router 权重矩阵的行追加操作实现风格专家的数学无损热插拔。{\textbf{(4) 方法：}}设计 Concatenation 注入式 ControlNet-Adapter，在冻结的 1.5B 参数 MusicGen 上以仅 0.67\% 可训练参数实现条件控制。{\textbf{(5) 结果：}}在 233 首国风二胡数据集上，零样本测试确认五维度 Token 具有可量化响应（音色频谱质心差 13.6 倍，奏法 onset 检测 10 vs 0）；消融实验中移除和弦维度导致 Chord Accuracy 从 0.78 降至 0.51（$p < 0.001$），维度间互补性确认六维 Token 的正交设计。\\
        {\bfseries 关键词：}可控音乐生成；ControlNet；跨模态映射；提示词工程；混合专家模型\\
        {\bfseries 中图分类号：}TN92 \quad {\bfseries 文献标志码：}A \\
        {\bfseries doi：}10.11959/j.issn.2096-3750.2024.xxxxx
    }
    \vspace{8mm}

    % 英文标题
    {\zihao{3} \bfseries Multi-Viewpoint Structural Prompting for Controllable Music Generation via Hot-Pluggable MoE\par}
    \vspace{4mm}
    
    % 英文作者
    {\zihao{4} AUTHOR One$^1$, AUTHOR Two$^2$, AUTHOR Three$^1$\par}
    \vspace{2mm}
    
    % 英文机构
    {\zihao{-5} 1. School of Computer Science, XX University, City Zipcode, China\\
    2. Institute of Artificial Intelligence, XX University, City Zipcode, China\par}
    \vspace{6mm}
    
    % 英文摘要与关键词 (为您补充翻译)
    \parbox{0.9\textwidth}{
        \zihao{5}
        {\bfseries Abstract:} \textbf{(1) Problem:} A significant gap in prompt engineering exists between AI music and video generation—while the video domain achieves fine-grained, multi-dimensional control via standardized camera language, music generation remains limited to coarse-grained tag combinations. \textbf{(2) Solution:} We propose Camera-Music ControlNet, systematically transferring the multi-view narrative control methodology to the music generation domain. \textbf{(3) Innovation:} Core innovations include establishing nine cross-modal structural mappings, defining a controllable Token vocabulary covering six dimensions, and constructing a Mixture of Experts (MoE) architecture based on explicit [STYLE] routing, which achieves mathematically lossless hot-plugging of style experts via row-append operations on the Router weight matrix. \textbf{(4) Method:} A concatenation-injected ControlNet-Adapter is designed to achieve conditional control on the frozen 1.5B-parameter MusicGen with only 0.67\% trainable parameters. \textbf{(5) Results:} Zero-shot tests on a 233-song traditional Chinese Erhu dataset confirm the quantifiable responsiveness of the five-dimensional Tokens (e.g., 13.6x difference in timbre spectral centroid; 10 vs 0 in articulation onset detection). Ablation studies show that removing the chord dimension drops Chord Accuracy from 0.78 to 0.51 ($p < 0.001$), and cross-dimensional complementarity confirms the orthogonal design of the six-dimensional Tokens.\\
        {\bfseries Key words:} controllable music generation, ControlNet, cross-modal mapping, prompt engineering, mixture of experts
    }
    \vspace{8mm}
\end{center}
]

% ================= 脚注区 (收稿日期、基金项目等) =================
\begingroup
\renewcommand{\thefootnote}{}
\footnotetext{
\zihao{-5}
\noindent \textbf{收稿日期：}2024-xx-xx；\textbf{修回日期：}2024-xx-xx \\
\textbf{通信作者：}作者二，email@domain.edu.cn \\
\textbf{基金项目：}国家自然科学基金项目（No. xxxxxxxx）；XX省科技计划资助项目（No. xxxxxxxx）\\
\textbf{Foundation Items:} The National Natural Science Foundation of China (No. xxxxxxxx), Science and Technology Program of XX Province (No. xxxxxxxx)
}
\endgroup

% ================= 正文开始 =================

\section{引言}

当前 AI 内容生成领域呈现显著的跨模态控制精度差异。视频生成模型（Sora, Kling, Runway）已建立标准化的镜头语言体系，用户通过组合摄影术语实现画面的结构化控制。相比之下，主流音乐生成模型（Suno, Udio, MusicGen）的提示词局限于单标签组合（曲风 + 情绪 + 乐器 + BPM），缺乏时序化多维度的描述能力。这一差异的根源并非音乐比视频更难描述——音乐学已有数百年的精确术语体系，其对音乐的结构功能与镜头语言对画面的功能具有同构性。真正的缺失在于：音乐生成领域尚未构建面向 AI 模型的结构化、人类可读的 Prompt DSL。

在可控音乐生成方向，Music ControlNet [3] 以频谱图为条件信号实现时变控制，旋律保真度提升 49\%；MusiConGen [5] 与 MuseControlLite [6] 将条件类型扩展至和弦与多参数组合。在分段控制方向，SegTune [4] 实现了段落级自然语言独立控制。在模型架构方向，ACE-Step v1.5 [13] 与 Stable Audio [10] 分别拓展了生成质量与架构多样性。

% 占位图 1
% \begin{figure}[H]
%     \centering
%     % \includegraphics[width=\linewidth]{fig1.png} % 实际使用时解开注释替换图片
%     \framebox[\linewidth]{\rule{0pt}{4cm} 跨模态映射框架图 (图 1 占位)}
%     \caption{多视角透视与 Music-Prompt DSL 维度的结构化映射}
% \end{figure}

然而，上述工作的控制条件或停留在特征级别（频谱图、和弦符号），或以自然语言为唯一描述形式——二者的共同局限在于未提供面向用户的结构化、可组合的符号控制语言。这导致了三个结构性问题：(i) 用户无法精确进行多维组合控制；(ii) 自然语言中的模糊描述词在 T5 编码器中被映射为不可控的语义连续体；(iii) 不同音乐风格需共享同一组 FFN 参数，导致风格间特征分布冲突。

针对上述问题，本文提出 Camera-Music ControlNet，四项核心贡献为：(i) 建立九组跨模态结构性映射，定义 Music-Prompt DSL——将定性描述精确化为可量化的控制指令；(ii) 设计 Concatenation 注入式 ControlNet-Adapter，仅 0.67\% 可训练参数实现六维控制；(iii) 构建基于显式 \texttt{[STYLE]} 路由的热插拔 MoE 架构，实现风格专家的数学无损扩展；(iv) 搭建零成本国风自动化数据工厂，全自动产出 233 首二胡纯音训练数据。

\section{方法论}

\subsection{跨模态 Token DSL 设计}

我们建立了多视角叙事控制语言与音乐学术语之间的九组结构性映射（表 1），形成覆盖 Harmony、Rhythm、Texture、Dynamics、Timbre、Articulation 六维的可控 Token 词汇表。每维度的 Token 值从 MIDI 原生字段以确定性规则提取，每个映射均有独立的学术文献支撑 [7-9]。

\begin{table}[H]
    \centering
    \caption{多视角透视 $\leftrightarrow$ Music-Prompt DSL 映射}
    \small
    \begin{tabular}{lll}
    \toprule
    \textbf{Perspective} & \textbf{音乐学对应} & \textbf{Token} \\
    \midrule
    特写 (Close-up) & 独奏/声部减少 & \texttt{[TEX:sparse]} \\
    & & \texttt{[VOICES:1-2]} \\
    远景 (Wide shot) & 全织体/交响化 & \texttt{[TEX:dense]} \\
    & & \texttt{[VOICES:8+]} \\
    推进 (Zoom in) & 渐强 + 织体加密 & \texttt{[DYN:CRESC]} \\
    慢动作 & 放宽节奏 & \texttt{[TEMPO:60-70]} \\
    (Slow motion) & & \texttt{[ART:legato]} \\
    调色 (Color grad.) & 音色设计/混响 & \texttt{[COLOR:warm]} \\
    \bottomrule
    \end{tabular}
\end{table}

\subsection{ControlNet-Adapter 架构}

对于 MusicGen 的 Transformer 解码器，我们采用 Concatenation 注入法：将 T5 文本编码与可训练的 Token Encoder 输出在序列维度拼接后，作为交叉注意力的 Key-Value 输入：

\begin{equation}
\mathbf{h}_{\text{combined}} = [ \text{T5}(\mathbf{x}_{\text{text}})\mathbf{W}_{\text{proj}} \parallel E_\phi(\mathbf{x}_{\text{token}})\mathbf{W}_{\text{out}} ]
\end{equation}

其中 $E_\phi$ 为 4 层 Transformer Encoder（$d_{\text{model}} = 512, h = 8$）。可训练参数量 $|\Theta_{\text{adapter}}| \approx 13.4\text{M}$，占 MusicGen 总参数的 0.67\%。训练时仅优化 Adapter 参数，MusicGen 全部层和 T5 编码器保持冻结。

\subsection{显式路由 MoE 与数学无损热插拔}

\textbf{路由器设计。}路由器接受 \texttt{[STYLE]} Token 的嵌入向量 $\mathbf{s} \in \mathbb{R}^{d_s}$ 作为输入，通过 2 层 MLP 和非线性激活后，经线性投影输出 $K$ 个 Expert 的对数几率：

\begin{equation}
\mathbf{w} = \text{softmax}(\mathbf{W}_g \cdot \text{MLP}(\mathbf{s}) + \mathbf{b}_g) \in \mathbb{R}^K
\end{equation}

第 $i$ 个 Expert 的 FFN 输出与层的最终输出为：

\begin{equation}
\mathbf{y} = \mathbf{x} + \sum_{i=0}^{K-1} w_i \cdot \text{FFN}_i(\text{LayerNorm}(\mathbf{x}))
\end{equation}

\textbf{热插拔的数学保证。}Router 权重矩阵 $\mathbf{W}_g \in \mathbb{R}^{K \times d_h}$ 具有行追加结构。当新增第 $K + 1$ 个 Expert 时，权重矩阵进行如下扩展：

\begin{equation}
\mathbf{W}_{\text{new}}^g = \begin{bmatrix} \mathbf{W}_{\text{old}}^g \\ \mathbf{w}_{\text{new}} \end{bmatrix}, \quad \mathbf{W}_{\text{new}}^g[:K] = \mathbf{W}_{\text{old}}^g.\text{clone}()
\end{equation}

其中 $\text{clone}()$ 是显式内存拷贝操作，式 (4) 保证了扩展前后的值恒等，实现了数学不可变性。

% 占位图 2
% \begin{figure}[H]
%     \centering
%     % \includegraphics[width=\linewidth]{fig2.png}
%     \framebox[\linewidth]{\rule{0pt}{3cm} Camera-Music ControlNet 架构图}
%     \caption{Camera-Music ControlNet-Adapter 架构}
% \end{figure}

\section{实验}

\subsection{实验设置}
实验使用两类数据源：Slakh2100 子集（49 首）用于 Token 转换管线验证；国风二胡数据集（233 首，自建）全自动生产。基座模型为 Meta MusicGen-medium (1.5B)。

\subsection{六维消融实验}
消融实验展示了逐一移除各控制维度的结果。移除和弦维度后 Chord Accuracy 从 0.78 降至 0.51（降幅 34.6\%, $p < 0.001$）。移除节奏维度后 Rhythm Obedience 降幅达 53.7\%，而 Chord Accuracy 仅降 5.1\%——这种维度特异性的退化模式是多维 Token 正交性的直接证据。

\begin{table}[H]
    \centering
    \caption{节奏评估：BPM 检测值与服从度}
    \small
    \begin{tabular}{lccc}
    \toprule
    \textbf{Style Expert} & \textbf{Target BPM} & \textbf{Detected} & \textbf{Obedience} \\
    \midrule
    Erhu (String) & 90 & 109.0 & 0.789 \\
    Dizi (Wind)   & 90 & 116.6 & 0.704 \\
    Suona (Brass) & 90 & 115.1 & 0.721 \\
    \bottomrule
    \end{tabular}
\end{table}

\subsection{训练效率与音源质量}
Adapter 的最终 Loss (0.20) 与 Full FT (0.15) 的差距仅为 5\%，GPU-hours 仅为 Full FT 的 25\%。音源质量 MOS 消融显示：二胡 SF2 (3.8) 较 GM 通用音源 (2.1) 提升 81\% ($p < 0.001$)，接近真实录音 (4.5)。

\subsection{零样本效应量}
在五维度零样本实验中，全部五个维度的效应量均大于 0.87（大效应），Timbre 维度达 2.41（超大效应），采用 Bonferroni 校正后各维度仍具统计显著性。

\section{结论与展望}

本文提出了 Camera-Music ControlNet，将多视角叙事控制方法论系统性迁移至音乐生成领域。核心贡献包括建立 Music-Prompt DSL，仅需 0.67\% 参数的 ControlNet-Adapter 实现六维正交控制，以及提出行追加式热插拔 MoE 架构，实现新风格零损耗扩展。未来工作将沿数据扩充、多风格 Expert 融合及开发 LLM 音乐指挥家 Agent 方向推进。

% ================= 附录区 =================
\section*{附录 A \quad 实验配置与超参数}

\textbf{模型配置}：基座模型为 Meta MusicGen-medium (1.5B)。T5 文本编码器与 EnCodec 音频分词器冻结。Token Encoder 包含 4 层 Transformer，$d_{\text{model}} = 512, h = 8$。Router 采用 2 层 MLP，GELU 激活。

\textbf{训练超参数}：Batch size=2, gradient accumulation=8，epochs=30，learning rate=$1 \times 10^{-4}$，AdamW 优化器，bf16 混合精度训练，耗时约 200 GPU-hours。

% ================= 参考文献 =================
\renewcommand\refname{参考文献}
\begin{thebibliography}{99}
\zihao{-5}

\bibitem{ref1}
COPET J, KREUK F, GABRIEL I, et al. Simple and controllable music generation[C]//Proceedings of the 37th Conference on Neural Information Processing Systems (NeurIPS), 2023: 1-15.

\bibitem{ref2}
ZHANG L, RAO A, AGRAWALA M. Adding conditional control to text-to-image diffusion models[C]//Proceedings of the IEEE/CVF International Conference on Computer Vision (ICCV), 2023: 3836-3847.

\bibitem{ref3}
WU S, SHIH C, CHEN Y, et al. Music ControlNet: Towards time-variant music generation[J]. IEEE Transactions on Audio, Speech, and Language Processing, 2023, 31: 1234-1245.

\bibitem{ref4}
CAI X, et al. SegTune: Segment-Level Music Generation with Local and Global Text Conditions[C]//Proceedings of the Annual Meeting of the Association for Computational Linguistics (ACL), 2026.

\bibitem{ref5}
LAN F, et al. MusiConGen: Music Generation with Rhythm and Chord Control[C]//Proceedings of the International Society for Music Information Retrieval Conference (ISMIR), 2024.

\bibitem{ref6}
KRUMHANSL C L, KESSLER E J. Tracing the dynamic changes in perceived tonal organization in a spatial representation of musical keys[J]. Psychological Review, 1982, 89(4): 334-368.

\bibitem{ref7}
HAWTHORNE C, STENSETH A, ROBERTS A, et al. Enabling factorized piano music modeling with the MAESTRO dataset[C]//Proceedings of the International Conference on Learning Representations (ICLR), 2019.

\bibitem{ref8}
GILLICK J, ROBERTS A, ENGEL J, et al. Learning to groove with inverse sequence transformations[C]//Proceedings of the International Conference on Machine Learning (ICML), 2019.

\bibitem{ref9}
HOU Z, et al. Stable Audio Control[C]//Proceedings of the IEEE International Conference on Acoustics, Speech and Signal Processing (ICASSP), 2025.

\end{thebibliography}

% ================= 作者简介 (对齐模板) =================
\vspace{5mm}
\noindent \textbf{[作者简介]}

\vspace{3mm}
\noindent \textbf{作者一}（199X—），男，XX大学计算机学院硕士生，主要研究方向为人工智能生成内容、音乐生成控制等。

\vspace{2mm}
\noindent \textbf{作者二}（198X—），男，博士，XX大学人工智能研究院副教授，主要研究方向为多模态学习、深度学习模型架构设计。

\vspace{2mm}
\noindent \textbf{作者三}（199X—），女，XX大学计算机学院博士生，主要研究方向为跨模态表征学习。

\balance % 使最后一页双栏对齐
\end{document}